\section{Comparative Evaluation}

We compare three models spanning the three architectural generations: a Vanilla VQA baseline, LXMERT, and BLIP.

\subsection{Experimental Setting}

\minisection{Vanilla VQA} We implemented this baseline from scratch following the architecture proposed by Antol et al.~\cite{antol2015vqa}, with a few deliberate changes over the original design. We replace their VGGNet image encoder with ResNet-101~\cite{he2016resnet} (pretrained on ImageNet, final classification layer removed), a stronger backbone unavailable at the time of the original paper; images are encoded offline and cached to disk as 2048-dimensional feature vectors, which also avoids loading the 44-million-parameter CNN during training, prohibitive on a single Kaggle T4 GPU. Questions are tokenized with NLTK and encoded by a 2-layer LSTM~\cite{hochreiter1997lstm} whose embedding layer is initialized with GloVe~\cite{pennington2014glove} rather than learned from scratch: pretrained embeddings give the model better semantic representations from the start and speed up convergence on limited GPU time; the embeddings are then fine-tuned during training.

The fusion mechanism is a Hadamard (element-wise) product between L2-normalized image and question representations, each first projected to 1024 dimensions through an MLP with LayerNorm and Tanh. A two-layer classifier with LayerNorm, ReLU, and Dropout produces logits over a fixed answer vocabulary.

Two key design choices deserve mention. First, the answer vocabularies (3129 for VQA v2, 1842 for GQA) were chosen to match figures commonly used in VQA literature: large enough to cover most questions, but not so large that many answers would have only a handful of training examples each, too few to learn reliably given our limited GPU budget. Second, the model uses different loss functions per dataset: binary cross-entropy with soft labels for VQA v2 (where answers are inherently ambiguous) and standard cross-entropy for GQA (where each question has one ground-truth answer). Training uses Adam with learning rate $5 \times 10^{-4}$, batch size 256, and early stopping.

\minisection{LXMERT} We evaluate two pretrained checkpoints from HuggingFace: \href{https://huggingface.co/unc-nlp/lxmert-vqa-uncased}{\texttt{unc-nlp/lxmert-vqa-uncased}} for VQA v2 and \href{https://huggingface.co/unc-nlp/lxmert-gqa-uncased}{\texttt{unc-nlp/lxmert-gqa-uncased}} for GQA. The model processes 36 pre-extracted Faster R-CNN region features per image (2048-dimensional) through three Transformer encoders. Like Vanilla VQA, LXMERT operates under closed-vocabulary classification: the output layer produces logits over a fixed answer set, and predictions are obtained via argmax. Inference is zero-shot: we use the provided classification heads without fine-tuning.

\minisection{BLIP} We use the single available checkpoint \href{https://huggingface.co/ybelkada/blip-vqa-base}{\texttt{ybelkada/blip-vqa-base}} from HuggingFace, loaded in \texttt{fp16} for memory efficiency. The model pairs a ViT-B/16 vision encoder with a BERT-based encoder-decoder, generating answers via greedy decoding (max 10 tokens). Output is free-form text rather than a vocabulary index. Since no GQA-specific BLIP checkpoint exists, GQA evaluation is zero-shot on the VQA v2 checkpoint.

\subsection{Evaluation Protocol}

For VQA v2, predictions on test-dev and test-standard are submitted to the Eval.AI server\footnote{Link to the VQA v2 challenge: \url{https://eval.ai/web/challenges/challenge-page/830/overview}}, which computes accuracy with per-type breakdowns. For GQA, the official competition is closed, so we evaluate directly on the test-dev balanced split by comparing predictions to ground-truth answers, reporting overall, binary (yes/no), and open accuracy.

\subsection{Results}

\begin{table}[t]
\centering
\small
\begin{tabular}{lcccc}
\toprule
\textbf{Model} & \textbf{Y/N} & \textbf{Num} & \textbf{Other} & \textbf{Overall} \\
\midrule
\multicolumn{5}{c}{\textit{VQA v2 test-dev}} \\
\midrule
Vanilla & 75.23 & 35.70 & 46.03 & 56.88 \\
LXMERT & 86.76 & 52.95 & 60.17 & 70.30 \\
BLIP   & 92.25 & 54.55 & 66.91 & 75.95 \\
\midrule
\multicolumn{5}{c}{\textit{VQA v2 test-standard}} \\
\midrule
Vanilla & 75.55 & 35.43 & 46.14 & 57.17 \\
LXMERT & 86.41 & 52.89 & 60.28 & 70.31 \\
BLIP   & 92.32 & 54.39 & 66.81 & 76.03 \\
\bottomrule
\end{tabular}
\caption{VQA v2 results (\% soft accuracy, averaged over 10 annotators).}
\label{tab:vqa_results}
\end{table}

\begin{table}[t]
\centering
\small
\begin{tabular}{lccc}
\toprule
\textbf{Model} & \textbf{Overall} & \textbf{Binary} & \textbf{Open} \\
\midrule
Vanilla & 47.49 & 64.46 & 37.95 \\
LXMERT  & 59.29 & 77.17 & 49.25 \\
BLIP    & 47.11 & 65.66 & 36.69 \\
\bottomrule
\end{tabular}
\caption{GQA test-dev balanced results (\% exact-match accuracy).}
\label{tab:gqa_results}
\end{table}

VQA v2 (Table~\ref{tab:vqa_results}) shows the expected generational progression: 56.88\% (Vanilla) $\rightarrow$ 70.30\% (LXMERT, +13.42) $\rightarrow$ 75.95\% (BLIP, +5.65). The largest gains appear between the Vanilla baseline and LXMERT, suggesting that cross-modal attention contributes more than the shift from classification to generation. All three models follow the same difficulty ranking: yes/no questions are easiest, counting is hardest.

GQA (Table~\ref{tab:gqa_results}) tells a different story. LXMERT leads at 59.29\%, while BLIP drops to 47.11\%, below even the Vanilla baseline (47.49\%). On the binary subset, LXMERT holds a 12-point advantage over BLIP (77.17\% vs.\ 65.66\%), and the gap widens on open questions (49.25\% vs.\ 36.69\%). Note that GQA's ``open'' category, unlike VQA v2's ``other'' column in Table~\ref{tab:vqa_results}, also includes counting questions, so the two are not directly comparable.

\subsection{Discussion}

The VQA v2 results are consistent with the architectural narrative. The Vanilla baseline's Hadamard fusion captures only first-order interactions between modalities; LXMERT's multi-layer cross-attention builds richer, context-dependent alignments; BLIP's generative pretraining adds a further boost, particularly on open-ended questions where its vocabulary is not artificially restricted. Our reproduced BLIP accuracy (75.95\%) is below the 78.25\% reported in the original paper~\cite{li2022blip}, plausibly due to two protocol differences rather than a flaw in our evaluation: the original paper fine-tunes and evaluates at $480\times480$ resolution, while our pipeline operates at $224\times224$, giving the model less visual detail; and the official BLIP codebase evaluates VQA by \emph{ranking} a fixed list of candidate answers, whereas our setup uses free-form greedy \emph{generation}, which is not constrained to match the exact wording of any ground-truth answer and is therefore more exposed to the exact-match penalty on semantically correct but differently phrased answers.

The GQA results demand a more nuanced reading. Unlike the modest VQA v2 gap discussed above, BLIP's much larger GQA shortfall cannot be attributed to evaluation protocol (rank vs.\ generation) or resolution alone; it stems primarily from zero-shot domain transfer, detailed next. LXMERT's advantage is not purely architectural, since it benefits from a checkpoint specifically fine-tuned on GQA. BLIP, lacking any GQA adaptation, must answer GQA questions using only its VQA v2 pretraining: on questions requiring compositional reasoning (e.g.,``What color is the car to the left of the tree?''), BLIP's free-form generation often misses the relational constraint, whereas LXMERT's fixed-vocabulary classification is constrained to answers the dataset expects. GQA's exact match further penalises open-vocabulary models for syntactic variation; BLIP is additionally disadvantaged by the absence of answer normalisation on its free-form output, unlike the fixed, canonicalised vocabularies of Vanilla and LXMERT. The expected gain from normalisation is small, however, and does not alter the ranking.

This reversal, with BLIP dominating VQA v2 but underperforming on GQA, illustrates a tension between generality and specialization. Generative models offer flexibility (answers beyond a fixed vocabulary) and strong in-domain performance, but their advantage shrinks when the evaluation distribution differs from pretraining. Closed-vocabulary classifiers trained on task-specific data remain competitive, particularly on benchmarks with controlled answer spaces.

A further limitation applies to all three models. Vanilla VQA and LXMERT are also capped by a fixed answer vocabulary: even if they reason correctly about the image, they can only output the closest available answer, so a wrong prediction does not necessarily mean an absence of reasoning. Still, the pattern is consistent with the well-known language-prior problem in VQA, where models partly guess from question wording rather than the image. BLIP is not restricted to a fixed list of answers, but its sharp drop on GQA suggests it relies on something similar: its answers still seem to depend on the patterns it picked up from VQA v2, since a model that truly reasoned about the image and question should not collapse so badly on a new dataset.

Counting remains a hard ceiling for all models. Even BLIP, the strongest system, reaches only 54.55\% on VQA v2 counting questions, compared to 92.25\% on yes/no. This gap reflects a persistent limitation of current architectures: detecting objects is necessary but not sufficient for counting them.
